% =============================================================================
%  text/THE ARGUMENT/ruler-dilemma.tex
%  Theory subsection: demonstrates an equation and marginal-returns logic.
%  理论小节示例：演示公式排版。
% =============================================================================

We formalize the ruler's decision as a comparison of marginal returns.
Let $R$ denote the ruler's probability of political survival and $S$ denote
the level of state capacity. The ruler allocates a fixed budget $B$ between
two activities: power consolidation $p$ and state-building $s$, with
$p+s=B$. The ruler solves

\begin{equation}
\max_{p,s} \; U(R(p),\, S(s)) \quad \text{subject to} \quad p+s=B,
\label{eq:ruler-problem}
\end{equation}

where $U$ is strictly increasing in both arguments.

The key comparative statics follow from the curvature of $R(\cdot)$.
When authority is insecure, $R$ is steep: each additional unit of
consolidation substantially reduces the probability of being overthrown, so
$\partial R/\partial p \gg \partial S/\partial s$ and the optimum shifts
toward fragmentation of elite networks. When authority is secure, $R$ is
flat, $\partial S/\partial s$ dominates, and the ruler invests in cohesive
elite networks instead.

\begin{equation}
\text{Strategy}^* =
\begin{cases}
\text{fragment elites}, & \text{if } \dfrac{\partial R}{\partial p}
  > \dfrac{\partial S}{\partial s};\\[2ex]
\text{foster elite cohesion}, & \text{if } \dfrac{\partial S}{\partial s}
  \geq \dfrac{\partial R}{\partial p}.
\end{cases}
\label{eq:strategy-choice}
\end{equation}

% 中文对照（可选）
% 式 (1) 刻画统治者在权力巩固与政权建设之间的资源配置问题；
% 式 (2) 给出边际收益比较下的最优策略选择。
This simple framework, summarized in equations~\eqref{eq:ruler-problem} and
\eqref{eq:strategy-choice}, generates the four strategies discussed in
Section~\ref{sec:four-strategies}.
